\section{Background of Study}

The proliferation of Internet of Things (IoT) sensor networks in smart building management has
created a pressing need for energy-aware operational strategies. A typical mid-sized commercial
building deployment may involve dozens of battery-powered or energy-harvesting sensor nodes
continuously reporting temperature, humidity, occupancy, and air-quality readings to a central
analytics platform. When these nodes sample and transmit data at fixed intervals regardless of
actual environmental variability or network conditions, substantial energy is wasted on redundant
or low-value transmissions.

Nexa Systems, a facilities-technology provider operating several smart building pilots in Kuala
Lumpur, reported that fixed-rate sensor polling accounted for approximately 34\% of total edge
device energy draw across their 2025 deployment portfolio. This observation motivates the central
question of this thesis: can a machine-learning-driven scheduler reduce energy consumption
without materially degrading the timeliness of environmental data delivered to downstream
analytics services?
